\chapter{Discussion}
\label{ch:discussion}

\section{What Running It Actually Taught Us}
\label{sec:defects}

Six defects were found by operating the system, not by reading the code.
They are reported in full because each one generalises beyond this
project.

\textbf{1. A healthcheck that proves liveness proves nothing about
work.} Restarting the Spark worker killed the archiver's executor. The
driver logged \code{Disconnected from Spark cluster!} and then simply sat
there: the query remained ``active'', so \code{awaitTermination()} never
returned. The container stayed up, \textbf{its healthcheck stayed green}
--- the metrics endpoint is served by the live driver --- and the archiver
silently stopped writing the master dataset for 26 minutes. The batch
layer then reconciled a day from an incomplete archive. The fix,
\code{streaming/watchdog.py}, monitors \emph{progress} rather than
liveness and exits so Docker restarts from the checkpoint. It was
verified by deliberately killing the executors: the speed layer reported
\code{vehicle\_state has not progressed for 132s}, exited, restarted and
resumed.

\textbf{2. A correct pipeline can produce a useless answer.} The first
full reconciliation flagged \textbf{35 of 50} vehicles as unprofitable,
with costs at 1.9$\times$ revenue. The pipeline was correct throughout;
the \emph{economics} were not calibrated --- fuel cost 7.35\,CU/km against
roughly 9.25\,CU revenue per driven kilometre, and service visits costing
3--6$\times$ a vehicle's daily revenue. After recalibration the figure is
\textbf{3--6 of 50} with a 38--48\% fleet margin. Only running it end to
end on realistic volumes surfaced this.

\textbf{3. Intervals chosen in real seconds are 60$\times$ longer in
business terms.} The odometer ledger checkpointed every 30 \emph{real}
seconds --- which at \code{COMPRESSION=60} is 30 \emph{simulated}
minutes. A container restart lost roughly 90\,km of ledger, and the batch
layer flagged \textbf{21 of 50} vehicles for a distance mismatch the
partner had not caused. A check that fires because of our own restart is
worse than no check: it trains the reader to ignore it. This is the
mismatch visible as the 21 distance mismatches on 2024-01-02 in
Table~\ref{tab:four-days}.

\textbf{4. Shared code publishes shared metrics.} The batch gauges were
defined in \code{common/metrics.py}, which every service imports --- so
all five services exported
\code{fleet\_batch\_last\_success\_timestamp} at its default of 0.
Prometheus saw five series, four permanently zero, and
\code{BatchNotRunRecently} fired forever on the bogus ones. A
permanently-firing false alert is worse than no alert.

\textbf{5. An alert that fires on a normal operation is noise.} The
expense SLA was evaluated against every file version, so a corrected
\code{v2} --- by definition sent after the deadline --- marked every
resubmission late. The SLA now governs first delivery only.

\textbf{6. Acceptance checks must actually be run.} The ``fresh clone
$\rightarrow$ one command'' criterion exposed two defects in the Windows
helper script that no amount of reading would have found: PowerShell
converting Docker's normal stderr progress into a terminating error, and
a wrapper function binding Docker's \code{-d} flag to itself, so the
stack came up in the foreground.

The common thread is that every one of these was invisible to code
review and visible within minutes of operating the system and checking
the numbers against expectations.

\section{Limitations and Threats to Validity}
\label{sec:limitations}

Stated plainly, because a report claiming no weaknesses is not believed.

\begin{itemize}
  \item \textbf{Simulated day 1 is partial} (06:00 start), so vehicles on
        night shifts work only part of it while carrying a full day's
        standing cost. Its 42.8\% margin is therefore not directly
        comparable with the other three days.
  \item \textbf{The revenue drift is usually zero} for the reasons
        measured in Section~\ref{sec:drift-zero}. The mechanism is
        visible in \code{fleet\_late\_rows\_dropped\_total}; its
        financial impact usually is not. A reader who wanted a dramatic
        number will not find one here.
  \item \textbf{\code{becoming\_unprofitable} requires three reconciled
        days}, so it cannot be demonstrated in a short session. This is
        correct behaviour --- we refuse to infer a trend from two points
        --- but it constrains the demo.
  \item \textbf{``Lemon'' vehicles are flagged more often, not always}
        --- 25\% of vehicle-days against 12.8\% for the rest of the fleet.
        A high-maintenance vehicle only looks bad on a day it breaks, and
        a low-demand vehicle earns less but also spends less. Realistic,
        but a demo should name a specific vehicle rather than promise
        that every lemon appears.
  \item \textbf{Batch Spark runs in local mode}, so the project
        demonstrates distributed \emph{streaming} but not distributed
        \emph{batch}.
  \item \textbf{Downtime creates genuine data gaps.} The simulated clock
        is derived from wall-clock time, so while the stack is stopped,
        simulated hours pass with no events emitted. A fifth simulated
        day was measurably degraded by our own deliberate outage testing.
        The pipeline reports the reduced revenue faithfully rather than
        fabricating it --- correct behaviour, but it means operational
        downtime is indistinguishable from low demand in the data.
  \item \textbf{One Kafka broker, RF~1.} Replication is out of scope, so
        broker loss means losing anything not yet archived.
  \item \textbf{Single-run measurements.} The four-day table is one run
        of one configuration. The drift figures in particular are small
        integer counts, so day-to-day variation is dominated by the
        Poisson-like arrival of late \code{trip\_end} events rather than
        by anything systematic. We report the arithmetic in
        Table~\ref{tab:drift-arithmetic} precisely so the reader can
        judge the expected variance rather than over-reading four
        samples.
\end{itemize}

\section{At Production Scale}
\label{sec:production}

Everything below is a deliberate simplification for a two-week,
single-laptop project, paired with what would actually be required if
this ran a real fleet.

\subsection{Ingestion}

\textbf{Today:} one Kafka broker, replication factor~1, 6 partitions,
7-day retention. \textbf{At scale:} a minimum of three brokers with RF~3
and \code{min.insync.replicas=2}. Today a broker loss means losing
anything not yet archived --- acceptable for a demo, indefensible for
revenue data. Partition count would be driven by measured consumer
throughput, but the \emph{keying} would not change: \code{vehicle\_id}
keying is a correctness constraint, not a tuning knob, and it caps useful
parallelism at the number of vehicles.

We version events with \code{schema\_version} and validate defensively. A
production deployment would put a \textbf{schema registry} (Avro or
Protobuf) in front of the topic so incompatible producers are rejected at
publish time rather than discovered in the DLQ. The DLQ would also gain
an automated replay path; today replay is manual.

\subsection{Storage}

\textbf{Today:} Parquet on a Docker volume, partitioned by
\code{sim\_date}. \textbf{At scale:} object storage --- already a
one-variable change, since every path is built from \code{LAKE\_ROOT}.
Two additions would matter quickly. A \textbf{table format} (Iceberg or
Delta Lake) would give atomic partition overwrites, schema evolution and
time travel; our \code{load\_batch\_views} achieves atomicity inside
PostgreSQL via delete-then-insert in a transaction, but the \emph{lake}
copy has no such guarantee, so a failed curated write can leave a partial
partition. \textbf{Compaction} would become mandatory: the archiver
writes every 30 seconds, a deliberate small-files trade-off that at real
volume would make the batch \emph{read} slower than the batch
\emph{compute}.

\subsection{Processing}

\textbf{Today:} Spark standalone, one worker, 3 cores, both streaming
applications co-resident. \textbf{At scale:} Kubernetes with the Spark
operator, giving per-application resource isolation and autoscaling. The
single most valuable change would be \textbf{decoupling the archiver from
the speed layer entirely} --- they already have separate checkpoints, but
they currently compete for the same worker's cores, which is how we lost
26 minutes of archiving (Section~\ref{sec:defects}).

Batch would stop running in local mode. That choice is sound at 12k
events/day and wrong at 12M; \code{compute\_vehicle\_day} is already
written as ordinary Spark SQL, so the change is a submission target, not
a rewrite.

\subsection{Serving}

\textbf{Today:} a single PostgreSQL instance serving both layers, queried
directly by Grafana. \textbf{At scale:} the access patterns genuinely
diverge. The \code{rt\_*} tables are high-frequency point lookups and
upserts; \code{batch\_*} is analytical scans over history. We would keep
PostgreSQL for the real-time tables, add read replicas for Grafana so
dashboard queries cannot slow the streaming sinks, and move
\code{batch\_vehicle\_daily} to a columnar warehouse once history
outgrew a single node. The API would gain caching on \code{/fleet/live}
--- it is read far more often than the underlying data changes --- and
authentication, which is explicitly out of scope here.

\subsection{Observability}

\textbf{Today:} Prometheus, Alertmanager, Grafana, structured JSON logs,
correlation IDs. \textbf{At scale:} three additions, in priority order.

\begin{enumerate}
  \item \textbf{Log aggregation} (Loki, or ELK). Today diagnosing a
        cross-service issue means \code{docker compose logs} per service.
        The logs are already structured JSON with a \code{stage} field
        specifically so they can be shipped and queried without
        reformatting --- this is a deployment step, not a code change.
  \item \textbf{Real distributed tracing.} As argued in
        Section~\ref{sec:tracing}, spans buy little for a
        Kafka-and-Parquet pipeline at this size. That calculus flips once
        there are multiple upstream producers and several consuming
        services, where ``which producer caused this DLQ spike?'' stops
        being answerable from correlation IDs alone.
  \item \textbf{SLOs with error budgets}, replacing fixed
        thresholds~\cite{beyer2016sre}. \code{IdleVehiclesHigh > 8} is a
        magic number tuned to a 50-vehicle fleet; it does not survive the
        fleet doubling. An SLO expressed as a \emph{ratio} would.
\end{enumerate}

\subsection{The Simulated Clock}

This is the largest single simplification and its consequence is worth
restating. Because simulated time is derived from wall-clock elapsed
time, \textbf{the pipeline being down does not pause the world} ---
simulated hours keep passing with no events emitted, creating a genuine
data gap. We observed exactly this: a simulated day degraded measurably
during our own outage testing.

The pipeline reports the reduced revenue faithfully rather than
fabricating it, which is the correct behaviour. But it means
\textbf{operational downtime is indistinguishable from low demand in the
data}. A production system would annotate known outage windows so that
downstream consumers --- and the \code{becoming\_unprofitable} trend rule
in particular --- can exclude them rather than mistaking an incident for a
commercial decline.
